%% Beginning of file 'sample7.tex'
%%
%% Version 7. Created January 2025.  
%%
%% AASTeX v7 calls the following external packages:
%% times, hyperref, ifthen, hyphens, longtable, xcolor, 
%% bookmarks, array, rotating, ulem, and lineno 
%%
%% RevTeX is no longer used in AASTeX v7.
%%
\documentclass[linenumbers,trackchanges]{aastex7}
%%
%% This initial command takes arguments that can be used to easily modify 
%% the output of the compiled manuscript. Any combination of arguments can be 
%% invoked like this:
%%
%% \documentclass[argument1,argument2,argument3,...]{aastex7}
%%
%% Six of the arguments are typestting options. They are:
%%
%%  twocolumn   : two text columns, 10 point font, single spaced article.
%%                This is the most compact and represent the final published
%%                derived PDF copy of the accepted manuscript from the publisher
%%  default     : one text column, 10 point font, single spaced (default).
%%  manuscript  : one text column, 12 point font, double spaced article.
%%  preprint    : one text column, 12 point font, single spaced article.  
%%  preprint2   : two text columns, 12 point font, single spaced article.
%%  modern      : a stylish, single text column, 12 point font, article with
%% 		  wider left and right margins. This uses the Daniel
%% 		  Foreman-Mackey and David Hogg design.
%%
%% Note that you can submit to the AAS Journals in any of these 6 styles.
%%
%% There are other optional arguments one can invoke to allow other stylistic
%% actions. The available options are:
%%
%%   astrosymb    : Loads Astrosymb font and define \astrocommands. 
%%   tighten      : Makes baselineskip slightly smaller, only works with 
%%                  the twocolumn substyle.
%%   times        : uses times font instead of the default.
%%   linenumbers  : turn on linenumbering. Note this is mandatory for AAS
%%                  Journal submissions and revisions.
%%   trackchanges : Shows added text in bold.
%%   longauthor   : Do not use the more compressed footnote style (default) for 
%%                  the author/collaboration/affiliations. Instead print all
%%                  affiliation information after each name. Creates a much 
%%                  longer author list but may be desirable for short 
%%                  author papers.
%% twocolappendix : make 2 column appendix.
%%   anonymous    : Do not show the authors, affiliations, acknowledgments,
%%                  and author contributions for dual anonymous review.
%%  resetfootnote : Reset footnotes to 1 in the body of the manuscript.
%%                  Useful when there are a lot of authors and affiliations
%%		    in the front matter.
%%   longbib      : Print article titles in the references. This option
%% 		    is mandatory for PSJ manuscripts.
%%
%% Since v6, AASTeX has included \hyperref support. While we have built in 
%% specific %% defaults into the classfile you can manually override them 
%% with the \hypersetup command. For example,
%%
%% \hypersetup{linkcolor=red,citecolor=green,filecolor=cyan,urlcolor=magenta}
%%
%% will change the color of the internal links to red, the links to the
%% bibliography to green, the file links to cyan, and the external links to
%% magenta. Additional information on \hyperref options can be found here:
%% https://www.tug.org/applications/hyperref/manual.html#x1-40003
%%
%% The "bookmarks" has been changed to "true" in hyperref
%% to improve the accessibility of the compiled pdf file.
%%
%% If you want to create your own macros, you can do so
%% using \newcommand. Your macros should appear before
%% the \begin{document} command.
%%
\newcommand{\vdag}{(v)^\dagger}
\newcommand\aastex{AAS\TeX}
\newcommand\latex{La\TeX}
%%%%%%%%%%%%%%%%%%%%%%%%%%%%%%%%%%%%%%%%%%%%%%%%%%%%%%%%%%%%%%%%%%%%%%%%%%%%%%%%
%%
%% The following section outlines numerous optional output that
%% can be displayed in the front matter or as running meta-data.
%%
%% Running header information. A short title on odd pages and 
%% short author list on even pages. Note that this
%% information may be modified in production.
%%\shorttitle{AASTeX v7 Sample article}
%%\shortauthors{The Terra Mater collaboration}
%%
%% Include dates for submitted, revised, and accepted.
%%\received{February 1, 2025}
%%\revised{March 1, 2025}
%%\accepted{\today}
%%
%% Indicate AAS Journal the manuscript was submitted to.
%%\submitjournal{PSJ}
%% Note that this command adds "Submitted to " the argument.
%%
%% You can add a light gray and diagonal water-mark to the first page 
%% with this command:
%% \watermark{text}
%% where "text", e.g. DRAFT, is the text to appear.  If the text is 
%% long you can control the water-mark size with:
%% \setwatermarkfontsize{dimension}
%% where dimension is any recognized LaTeX dimension, e.g. pt, in, etc.
%%%%%%%%%%%%%%%%%%%%%%%%%%%%%%%%%%%%%%%%%%%%%%%%%%%%%%%%%%%%%%%%%%%%%%%%%%%%%%%%
%%
%% Use this command to indicate a subdirectory where figures are located.
%%\graphicspath{{./}{figures/}}
%% This is the end of the preamble.  Indicate the beginning of the
%% manuscript itself with \begin{document}.

\begin{document}

\title{Understanding the Contributions of the Milky Way and M31 to the Merger Remnant Shape (Research Assignment 4)}

%% A significant change from AASTeX v6+ is in the author blocks. Now an email
%% address is required for each author. This means that each author requires
%% at least one of the following:
%%
%% \author
%% \affiliation
%% \email
%%
%% If these three commands are not available for each author, the latex
%% compiler will issue an error and if you force the latex compiler to continue,
%% it will generate an incomplete pdf.
%%
%% Multiple \affiliation commands are allowed and authors can also include
%% an optional \altaffiliation to indicate a status, i.e. Hubble Fellow. 
%% while affiliations are indexed as footnotes, altaffiliations are noted with
%% with a non-numeric footnote that is set away from the numeric \affiliation 
%% footnotes. NOTE that if an \altaffiliation command is used it must 
%% come BEFORE the \affiliation call, right after the \author command, in 
%% order to place the footnotes in the proper location. Because non-numeric
%% symbols are used, \altaffiliation should be used sparingly.
%%
%% In v7 the \author command takes an optional argument which provides 
%% additional metadata about the author. Authors can provide the 16 digit 
%% ORCID, the surname (family or last) name, the given (first or fore-) name, 
%% and a name suffix, e.g. "Jr.". The syntax is:
%%
%% \author[orcid=0000-0002-9072-1121,gname=Gregory,sname=Schwarz]{Greg Schwarz}
%%
%% This name metadata in not shown, it is only for parsing by the peer review
%% system so authors can be more easily identified. This name information will
%% also be sent to the publisher so they can include it in the CROSSREF 
%% metadata. Including an orcid will hyperlink the author name to the 
%% author's ORCID page. Note that  during compilation, LaTeX will do some 
%% limited checking of the format of the ID to make sure it is valid. If 
%% the "orcid-ID.png" image file is  present or in the LaTeX pathway, the 
%% ORCID icon will appear next to the authors name.
%%
%% Even though emails are now required for each author, the \email does not
%% produce output in the compiled manuscript unless the optional "show" command
%% is used. For example,
%%
%% \email[show]{greg.schwarz@aas.org}
%%
%% All "shown" emails are show in the bottom left of the first page. Due to
%% space constraints, only a few emails should be shown. 
%%
%% To identify a corresponding author, use the \correspondingauthor command.
%% The command appends "Corresponding Author: " to the argument it appears at
%% the bottom left of the first page like the output from \email. 

\author[orcid=0000-0000-0000-0001,sname='Catherine Snyder']{Catherine Snyder}
\altaffiliation{Steward Observatory}
\affiliation{University of Arizona}
\email[show]{catitude19@arizona.edu}  

\collaboration{all}{Submitted April 11th, 2025}





\keywords{\uat{Local Group}{1} --- \uat{Dark Matter Halo}{1} --- \uat{Halo Shape}{1} --- \uat{Halo Spin}{1} --- \uat{Oblate/Prolate/Triaxial}{1}}



%% Use the \collaboration command to identify collaborations. This command
%% takes an optional argument that is either a number or the word "all"
%% which tells the compiler how many of the authors above the command to
%% show. For example "\collaboration[all]{(DELVE Collaboration)}" wil include
%% all the authors above this command.
%%
%% Mark off the abstract in the ``abstract'' environment. 


%% Keywords should appear after the \end{abstract} command. 
%% The AAS Journals now uses Unified Astronomy Thesaurus (UAT) concepts:
%% https://astrothesaurus.org
%% You will be asked to selected these concepts during the submission process
%% but this old "keyword" functionality is maintained in case authors want
%% to include these concepts in their preprints.
%%
%% You can use the \uat command to link your UAT concepts back its source.


%% From the front matter, we move on to the body of the paper.
%% Sections are demarcated by \section and \subsection, respectively.
%% Observe the use of the LaTeX \label
%% command after the \subsection to give a symbolic KEY to the
%% subsection for cross-referencing in a \ref command.
%% You can use LaTeX's \ref and \label commands to keep track of
%% cross-references to sections, equations, tables, and figures.
%% That way, if you change the order of any elements, LaTeX will
%% automatically renumber them.

\section{Introduction} 

The area of galaxy evolution this research will investigate is the impact galaxy mergers have on dark matter halos. This research will specifically look at mergers in the \textbf{Local Group} which is the gravitationally bound systems of galaxies that the Milky Way is a part of. Dark matter is believed to be dispersed throughout the universe in an almost web-like structure with a vast system of filaments. \textbf{Dark matter halos} are the areas in the galaxy where there is a high density of dark matter. The overall shape these halos take is the \textbf{Halo Shape}, and the rotation of the halo is the \textbf{Halo Spin}. The Halo Shape is \textbf{Triaxial}, which means it is described by its three axes. If one of the axes is proportionally smaller or larger, the shape of the halo can be \textbf{Prolate}, or \textbf{Oblate} ellipsoids. These halos are theorized to come about when multiple filaments intersect with one another. These dark matter halos are significant because they are where galaxies are able to form \cite{Frenk_2012}. 

A \textbf{Galaxy} is a system of stars that are gravitationally bound and have physical properties that cannot be explained by just the stars alone, we define galaxies as being made of stars and dark matter. Because these dark matter halos are the basic unit of cosmological structure, understanding dark matter halos is intrinsic to understanding cosmology. Researching the structure of dark matter halos, as well as halo formation and halo growth can give key insights into \textbf{Galaxy Evolution}, which is the formation and life cycle of a galaxy. Understanding halo growth requires understanding how these structures are impacted when galaxies merge. 

The main understanding of dark matter halo mergers comes from cosmological simulations of mergers set in near perfect settings. So far the general understanding finds that the size  and spin of the halo remnant are dependent on the energy of the orbit and the angular momentum of the orbit respectively. The shape of the remnants is primarily dependent on the initial separation between the two halos as well as the spin of the mergers \cite{2019MNRAS.487..993D}. as seen in Figure 1. Although the relationship between energy of a system and the spread (or how compacted/extended) of a remnant appears to be direct, the correlation between the energy of a system and the radius of the scaled density profile of the remnant appears to be an inverse relationship \cite{2019MNRAS.487.1008D}.

%% The "ht!" tells LaTeX to put the figure "here" first, at the "top" next
%% and to override the normal way of calculating a float position.
%% The asterisk after "figure" tells the compiler to span multiple columns
%% if a two column style is selected.
\begin{figure*}[ht!]
\plotone{Figure1.png}
\caption{Figure 10 from Drakos + MNRAS, 2019a, 487, 993. This figure shows the results from the paper's merger simulations. The top graphs are the particles of the remnants, the middle graphs are the density profiles, and the bottom graphs are the mass profiles. These results show the relations between radial orbits and prolate shapes, and tagential orbits and oblate shapes. This figure confirms the idea that the shape of the remnant is dependent on the spin of the mergers.
\label{fig:general}}
\end{figure*}



A lot of the understanding of dark matter halo mergers comes from simulations done with equal mass binary systems, which is a simplified setting for mergers, as it does not take into account mergers between different sized halos or mergers between more than two halos. These simulations do not take into account halo substructures \cite{2019MNRAS.487..993D}. One of the largest issues still being examined from these simulations is what causes the central density of merger remnants to decrease, because the current simulations see an increase in central density that does not match the drop seen in central density of halos as they grow \cite{2019MNRAS.487.1008D}. Other questions left open include how the relationships and general findings of these merger simulations might change when including more complicated and more common types of mergers, such as mergers between more than two halos. 


\section{Project} \label{sec:style}

In this paper, we will examine the change in the shape of the remnant during the collision, as well examining how the contributions from the two original galaxies disperse during and after the merger. 

This project will address the question of if mergers in complicated systems with more than one galaxy influencing the merger, will coincide with the findings from the simulations of simplified mergers. It will also look into how the particles from a specific galaxy diffuse through the remnant after the merger. 

Understanding what factor influence the shape of the remnants is important to galaxy evolution because it helps us understand how the galaxies in our universe evolve over time. Understanding merger remnant shapes can help to understand the future of younger galaxies, as well as give us insights into older galaxies. Overall, understanding the shape of the merger gives us insight into the life cycles of galaxies.

\section{Methodology} \label{sec:style}

The simulations used in this project are from the work done on interactions between the Milky Way galaxy and the LMC and SMC.\cite{2012MNRAS.421.2109B} These simulations are done using an N-body simulation, which is a simulation of a system of particles and the forces acting on and between the particles in that system. The N-body system used in Marel, Besla +2012 is modeled with three different types of particles from the stellar bulge, stellar disk, and dark matter halo. This simulation runs approximately 12 Gyr.

The particle types this research will be using for analysis are the Type 1 dark matter or halo particles. This research will be looking at lower resolution(VLowRes) snapshots after the two galaxies merge, so in roughly 6.5 Gyr, which correlates to snapshots 435 and beyond. This research will compare multiple snapshots after merger, including snapshots 435, 535, 635, and 735 to compare how the shape of the remnant as well as the dispersion of the particles from the two original galaxies, changes over time. 

The way the questions of this research will be addressed is by creating dark matter distribution histograms in the XY, YZ, and XZ dimensions at the different snapshots. The three histograms for each snapshot will have ellipses fitted to them that will give the axes for the shape of the remnant at that time. Creating ratios of these axes will determine if the remnant is oblate or prolate. Computing these ratios for each of the snapshots being used will provide us the shape of the remnant over time. 

There are several sets of plots being made for this project. The histograms used for the triaxiality mentioned in the paragraph above will is one of the sets of plots being created. The ratio of the axes from these histograms will also be plotted over time to show how the shape of the remnant changes over time. The other plot being made will look at the particle contributions from each of the original galaxies over time. The way this will be done is making another histogram of the remnant that looks at the ratio of Milky Way particles to M31 particles. The histogram of this ratio will show visually the areas domnated by Milky Way particles and the areas dominated by M31 particles in the remnant. Making this histogram for the different snapshots can show how the contributions from each galaxy distribute over time after the merger. 
The models below represent the questions being investigated by this project:
\begin{figure*}[ht!]
\plotone{Figure 2.1.png}
\caption{The two figures at the top are three dimensional plots of a sample of the particles before the merger(left) and after the merger(right). These visualize the investigation this project is doing into the shape of the remnant. The figure at the bottom is a histogram of the ratio of M31 to Milky Way particles. This visually represents the project's examination of the contributions from the two original galaxies to the final remnant.
\label{fig:general}}
\end{figure*}

Due to the path of collision the Milky Way and M31 are on, the expected outcome of this research is to find a 3D simulation with the qualities similar to the radial orbit remnants shown in Figure 10 from Drakos et al. The anticipated end result will have particles from each of the galaxies intermixing, but with one galaxy’s particle type dominating the center of the remnant while the other galaxy's particle types become more prominent the further you get from the center of the remnant. 

%% For this sample we use BibTeX plus aasjournalv7.bst to generate the
%% the bibliography. The sample7.bib file was populated from ADS. To
%% get the citations to show in the compiled file do the following:
%%
%% pdflatex sample7.tex
%% bibtext sample7
%% pdflatex sample7.tex
%% pdflatex sample7.tex

\bibliography{Resources}{}
\bibliographystyle{aasjournalv7}

%% This command is needed to show the entire author+affiliation list when
%% the collaboration and author truncation commands are used.  It has to
%% go at the end of the manuscript.
%\allauthors

%% Include this line if you are using the \added, \replaced, \deleted
%% commands to see a summary list of all changes at the end of the article.
%\listofchanges

\end{document}

% End of file `sample7.tex'.
